\documentclass[12pt]{exam}
\usepackage[margin=1in]{geometry}
\usepackage{amsmath}
\usepackage{amsfonts}
\usepackage{tikz}
\usepackage{tkz-euclide}
\usepackage{multicol}
\usepackage{tasks}
\usepackage{xcolor}
\usepackage{makecell}

\begin{document}

\begin{center}
\fbox{\fbox{\parbox{5.5in}{\centering Pre-Calculus 11 - Reasoning and Analyzing Units 1+2 (V1)}}}
\\
\vspace{5mm}
\makebox[0.75\textwidth]{Name :\enspace\hrulefill}

\vspace{3mm}
\makebox[0.75\textwidth]{Teacher's name:\enspace\hrulefill}
\end{center}
\begin{center}
    \begin{tabular}{ |wc{5.5cm}||wc{0.7cm}|wc{0.7cm}|wc{0.7cm}|wc{0.7cm}|wc{0.7cm}|wc{0.7cm}|wc{0.7cm}|wc{0.7cm}|}
    \hline
    \centering \textbf{Outcome} & 0 & 1 & 2 & 3 & 4 & 5 & 6& 7\\ 
    \hline
    \centering Reasoning and Analyzing& & & & & & & & \\ 
    \hline
    \end{tabular}
\end{center}

%\printanswers
\begin{questions}
\question For each of the following statements, identify if the statement is always true, sometimes true, or never true. Provide a brief justification.
\begin{parts}
    

\part Let $a,b,c$ be non-zero integers, then $ax^2+bx+c$ can be factored, if there are two integers $m,n$ such that $m\times n=ac$ and $m+n=b$.
\begin{solution}
    Always true. That's the correct determining factor.
\end{solution}
\vspace{\stretch{1}}
\part Let $A,B$ be any real number, then $A+B\sqrt{3}$ is an irrational number.
\begin{solution}
    Sometimes true. If $B=0$ or $B=\sqrt{3}$, then we can have a rational number if $A$ is rational. If $B$ does not remove the irrational component of $\sqrt{3}$, then this statement is true.
\end{solution}
\vspace{\stretch{1}}
\part If $A,B\geq 0$, then $A\sqrt{B}=\sqrt{A^2B}$.
\begin{solution}
    Always true. Method for entire radicals to mixed and vice versa.
\end{solution}
\vspace{\stretch{1}}
\part If $x^2+bx+c$ can be factored, then $x^2+bx+c=0$ has two solutions for $x$.
\begin{solution}
    Sometimes true. It can be one solution but if someone argues the one solution is a \textit{double}, ie. two of the same answer, then allow it.
\end{solution}
\vspace{\stretch{1}}
\part If $A>0$, then $\sqrt{A}>\sqrt[3]{A}$.
\begin{solution}
    Sometimes true. For $0<A<1$, the cubic root is larger.
\end{solution}
\vspace{\stretch{1}}
\end{parts}
\pagebreak

\question A student was asked to solve the equation and leave answers in exact form. Which step did he first make an error? Explain the error.
    \begin{flalign*}
         && (2x+1)^2&=4x-8 &&\\   
         && 4x^2+4x+1&=4x-8 &&{\textbf{Step 1}}\\
         && 4x^2+9&=0 && \textnormal{\textbf{Step 2}}\\
         && (2x+3)(2x-3)&=0 && \textnormal{\textbf{Step 3}}\\
         && x&=\pm\frac{3}{2} && \textnormal{\textbf{Step 4}}
    \end{flalign*}
\vspace{30mm}
    \question In each of the following, one item doesn't belong....
        \begin{tasks}(4)
            \task* Which expression doesn't belong? Explain.
        
                    \task[i. ] $\sqrt{8x}$ \task[ii. ] $\sqrt[3]{8x}$  \task[iii. ] $\sqrt{18x}$  \task[iv. ] $\sqrt{12x}$
        
            \vspace{45mm}
            \task* Which equation doesn't belong? Explain.
                \task*(2)[i. ] $x^2-13x+42=0$ \task*(2)[ii. ]$\sqrt{x-6}=x-6$ 
                \task*(2)[iii. ] $x^2+7x+6=0$ \task*(2)[iv. ] $(7-x)(6-x)=0$ 
            \vspace{40mm}
        \end{tasks}
\end{questions}
\end{document}